% Template LaTeX untuk Ujian Kualifikasi
% Program Studi Doktor Informatika
% Universitas Telkom

% Document Class
\documentclass[12pt,a4paper]{article}

% Packages
\usepackage[utf8]{inputenc}
\usepackage[T1]{fontenc}
\usepackage{mathptmx} % Times New Roman
\usepackage{courier}  % Courier for monospace
\usepackage{helvet}   % Helvetica for sans-serif
\usepackage{amsmath}
\usepackage[english,indonesian]{babel}
\usepackage{geometry}
\usepackage{setspace}
\usepackage{titlesec}
\usepackage{graphicx}
\usepackage{hyperref}
\usepackage{booktabs}
\usepackage{array}
\usepackage[format=hang,font=footnotesize,labelfont=bf]{caption}
\usepackage{enumitem}
\usepackage{parskip}
\usepackage{fancyhdr}
\usepackage{tocloft}
\usepackage{tabularx}
\usepackage{longtable}
\usepackage{ragged2e}  % For better text justification in bibliography
\usepackage{fancyvrb}

% Page Setup
\geometry{
    top=3cm,
    bottom=3cm,
    left=4cm,
    right=3cm
}

% Title Formatting
\titleformat{\section}
    {\normalfont\fontsize{16}{19}\bfseries\centering}{\thesection}{1em}{}
    
\titleformat{\subsection}
    {\normalfont\fontsize{12}{14}\bfseries}{\thesubsection}{1em}{}

\titleformat{\subsubsection}
    {\normalfont\fontsize{12}{14}\normalfont}{\thesubsubsection}{1em}{}

% Spacing
\onehalfspacing

% Header/Footer
\pagestyle{fancy}
\fancyhf{}
\fancyfoot[C]{\thepage}
\renewcommand{\headrulewidth}{0pt}

% Hyperref setup
\hypersetup{
    colorlinks=true,
    linkcolor=black,
    citecolor=black,
    filecolor=black,
    urlcolor=black,
}

% Table of Contents customization
\renewcommand{\contentsname}{DAFTAR ISI}
\renewcommand{\cfttoctitlefont}{\hfill\fontsize{16}{19}\selectfont\bfseries}
\renewcommand{\cftaftertoctitle}{\hfill}

% List of Tables customization
\renewcommand{\listtablename}{DAFTAR TABEL}
\renewcommand{\cftlottitlefont}{\hfill\fontsize{16}{19}\selectfont\bfseries}
\renewcommand{\cftafterlottitle}{\hfill}

% List of Figures customization
\renewcommand{\listfigurename}{DAFTAR GAMBAR}
\renewcommand{\cftloftitlefont}{\hfill\fontsize{16}{19}\selectfont\bfseries}
\renewcommand{\cftafterloftitle}{\hfill}

% List of Equations setup
\newcommand{\listequationsname}{DAFTAR PSEUDOCODE}
\newlistof{myequations}{equ}{\listequationsname}
\newcommand{\myequations}[1]{%
\addcontentsline{equ}{myequations}{\protect\numberline{\theequation}#1}\par}
\renewcommand{\cftequtitlefont}{\hfill\fontsize{16}{19}\selectfont\bfseries}
\renewcommand{\cftafterequtitle}{\hfill}
\setlength{\cftmyequationsindent}{0pt}

% Redefine \listofequations to use the custom list
\newcommand{\listofequations}{\listofmyequations}

% Subsection numbering format
\renewcommand{\thesubsection}{\arabic{section}.\arabic{subsection}}

% Custom Commands
\newcommand{\myTitle}[1]{\def\@myTitle{#1}}
\newcommand{\myName}[1]{\def\@myName{#1}}
\newcommand{\myNIM}[1]{\def\@myNIM{#1}}
\newcommand{\myYear}[1]{\def\@myYear{#1}}

% Redefine title
\makeatletter
% Initialize default values
\def\@myTitle{Judul Penelitian}
\def\@myNIM{303012510004}
\def\@myName{Rolly Maulana Awangga}
\def\@myYear{Tahun}

\renewcommand{\maketitle}{
    \begin{titlepage}
        \centering
        \vspace*{0.8cm}
        
        {\fontsize{14}{18}\selectfont
        \textbf{LAPORAN TUGAS BESAR PENERAPAN ALGORITMA}\\
        \textbf{PADA APLIKASI PASARKITA}\par}
        
        \vspace{0.9cm}
        
        {\fontsize{12}{14}\selectfont
        \textbf{MATA KULIAH STRATEGI ALGORITMA}\par}
        
        \vspace{1.1cm}
        
        \IfFileExists{logo.png}{
            \includegraphics[width=0.36\textwidth]{logo.png}
        }{
            {\fontsize{16}{19}\bfseries Logo\par}
        }
        
        \vspace{1.2cm}

        {\fontsize{12}{14}\selectfont
        \textbf{DISUSUN OLEH:}\par}

        \vspace{0.5cm}

        \begin{tabular}{@{} l l @{}}
        \textbf{RADITYA RIZKI RAHARJA} & \textbf{714240041} \\
        \textbf{RIDWAN HAKIM RAMADHAN} & \textbf{714240050} \\
        \textbf{MUHAMMAD RASHID AL SAVERO} & \textbf{714240006} \\
        \end{tabular}

        \vspace{0.9cm}

        {\fontsize{12}{14}\selectfont
        \textbf{DOSEN PENGAMPU:}\par}

        \vspace{0.5cm}

        {\fontsize{12}{14}\selectfont
        \textbf{RONI HABIBI, S.KOM., M.T., SFPC}\par}

        \vspace{1.0cm}
        
        {\fontsize{12}{15}\selectfont
        \textbf{PROGRAM STUDI DIV TEKNIK INFORMATIKA}\par
        \textbf{UNIVERSITAS LOGISTIK \& BISNIS INTERNASIONAL}\par
        \textbf{BANDUNG}\par
        \textbf{2026}\par}
        
    \end{titlepage}
    \newpage
}
\makeatother

% SOTA Table Environment
\newenvironment{sototable}
    {\begin{table}[htbp]
    \centering
    \caption{State-of-the-Art Penelitian}
    \label{tab:sota}}
    {\end{table}}

% Word Count Environments
\newenvironment{summary}
    {}
    {}

\newenvironment{background}
    {}
    {}

\newenvironment{literaturereview}
    {}
    {}

\newenvironment{researchgap}
    {}
    {}

\newenvironment{novelty}
    {}
    {}

% Bibliography Style - IEEE with BibTeX
\usepackage[
    backend=bibtex,   % Menggunakan BibTeX (biber tidak berfungsi di sistem ini)
    style=ieee,     % IEEE numeric citation style
    sorting=none,     % IEEE: sort by citation order
    maxbibnames=99,   % Tampilkan semua nama penulis
    dashed=false,     % Jangan gunakan dash untuk author yang sama
    url=false,        % Jangan tampilkan URL
    doi=true,         % Tampilkan DOI
]{biblatex}

\urlstyle{same} % Samakan font URL/DOI dengan teks utama

% BibTeX cannot read spaces in database paths; this file mirrors src/My Collection.bib.
\addbibresource{src/MyCollection.bib}
%\addbibresource{references.bib}
%\addbibresource{include.bib}
%\addbibresource{visualdecoding.bib} % Uncomment jika ingin menggunakan file bib tambahan

% Make \cite behave like \parencite (with parentheses)
\let\cite\parencite

% Custom citation formatting
\renewcommand*{\finalnamedelim}{\addspace\&\addspace}  % Ubah "and" jadi "&"
\renewcommand*{\nameyeardelim}{\addcomma\space}        % Tambahkan koma setelah author

% Remove "In:" from bibliography
\renewbibmacro{in:}{}

% Ensure pages are not bold
\DeclareFieldFormat{pages}{\normalfont #1}

% Custom format untuk memastikan ada spasi setelah "in"
\DeclareFieldFormat{booktitle}{\addspace\textit{#1}}
\DeclareFieldFormat{journaltitle}{\addspace\textit{#1}}

% Fix URL breaking untuk mencegah overflow
\setcounter{biburllcpenalty}{7000}
\setcounter{biburlucpenalty}{8000}

% Set bibliography to use ragged right alignment untuk mencegah overfull
\emergencystretch=1em
\appto{\bibsetup}{\sloppy\RaggedRight}



% Set document metadata
\myTitle{Laporan Analisis Penerapan Algoritma pada PasarKita}
\myNIM{303012510004}
\myName{Rolly Maulana Awangga}
\myYear{2026}

% Beginning of Document
\begin{document}

% Force TOC title to be uppercase (fix for babel override)
\renewcommand{\contentsname}{DAFTAR ISI}
\renewcommand{\listtablename}{DAFTAR TABEL}
\renewcommand{\listfigurename}{DAFTAR GAMBAR}

% Use Times New Roman as default font
\rmfamily

% Title Page
\maketitle

% Halaman Depan menggunakan angka Romawi
\pagenumbering{roman}

% DAFTAR ISI
\phantomsection
\addcontentsline{toc}{section}{DAFTAR ISI}
\tableofcontents
\clearpage

% DAFTAR TABEL
\phantomsection
\addcontentsline{toc}{section}{DAFTAR TABEL}
\listoftables
\clearpage

% DAFTAR PSEUDOCODE
\phantomsection
\addcontentsline{toc}{section}{DAFTAR PSEUDOCODE}
\listofequations
\clearpage

% RINGKASAN
\clearpage
\section*{RINGKASAN}
\phantomsection
\addcontentsline{toc}{section}{RINGKASAN}
\section*{RINGKASAN}
\addcontentsline{toc}{section}{RINGKASAN}

Perkembangan teknologi \textit{Artificial Intelligence} (AI) dan \textit{Machine Learning} (ML) menuntut penguatan kompetensi baru dalam pendidikan vokasional, khususnya bidang Rekayasa Perangkat Lunak. Siswa SMK tidak hanya perlu menguasai dasar pemrograman, tetapi juga membutuhkan pengalaman praktis membangun solusi berbasis data yang relevan dengan kebutuhan industri digital. Mitra kegiatan Pengabdian kepada Masyarakat ini adalah SMKN 2 Cimahi, Jurusan Rekayasa Perangkat Lunak, dengan permasalahan utama berupa terbatasnya akses pelatihan praktis \textit{Machine Learning} dan \textit{Computer Vision}, belum adanya pengalaman pembelajaran berbasis proyek kecerdasan buatan secara, serta belum optimalnya portofolio siswa untuk mendukung kesiapan kerja maupun studi lanjut.

Untuk menjawab permasalahan tersebut, kegiatan ini menyelenggarakan Workshop \textit{Machine Learning} untuk Prediksi Awal Kesehatan Tanaman melalui Analisis Citra Daun. Studi kasus dipilih karena bersifat visual, menggunakan \textit{dataset} publik gratis, tidak memerlukan perangkat keras khusus, dan relevan dengan kehidupan sehari-hari. Sistem diposisikan sebagai alat \textit{screening} awal, bukan diagnosis final, agar siswa memahami keterbatasan model dan pentingnya interpretasi hasil secara kritis. Pelaksanaan menggunakan pendekatan \textit{Project-Based Learning} yang mencakup penyusunan modul pelatihan, pengenalan konsep kecerdasan buatan, eksplorasi \textit{dataset} citra daun, pelatihan model klasifikasi menggunakan \textit{Google Colab}, evaluasi model, hingga pengembangan aplikasi prediksi sederhana berbasis \textit{Streamlit}. Evaluasi kegiatan dilakukan melalui \textit{pre-test} dan \textit{post-test}, penilaian \textit{mini project} kelompok, serta kuesioner kepuasan peserta, dengan dukungan mahasiswa ULBI sebagai fasilitator teknis.

Kegiatan ini diharapkan meningkatkan pemahaman dan keterampilan siswa dalam menerapkan \textit{Machine Learning} dan \textit{Computer Vision} secara praktis, menghasilkan bahan ajar yang dapat digunakan kembali oleh mitra, serta membentuk pengalaman belajar berbasis proyek yang memperkuat kompetensi digital siswa SMK.

\noindent\textbf{Kata Kunci} : \textit{Machine Learning}, \textit{Computer Vision}, \textit{Project-Based Learning}.


\clearpage

% BAB 1 - PENDAHULUAN ANALISIS
\setcounter{section}{0}
\pagenumbering{arabic} % Kembali ke angka Arab mulai dari Bab 1
\setcounter{page}{1}   % Reset halaman ke 1
\section*{BAB 1 \\ PENDAHULUAN ANALISIS}
\phantomsection
\addcontentsline{toc}{section}{BAB 1 PENDAHULUAN ANALISIS}
\refstepcounter{section}

\subsection{Konteks Aplikasi}

PasarKita adalah aplikasi marketplace digital untuk produk UMKM. Aplikasi ini mempertemukan tiga aktor utama: pembeli, penjual, dan admin. Pembeli menggunakan aplikasi untuk melihat katalog, mencari produk, memasukkan produk ke keranjang, melakukan checkout, dan melihat status pesanan. Penjual menggunakan aplikasi untuk mengelola produk, stok, pesanan, serta dashboard toko. Admin menggunakan aplikasi untuk memantau user, produk, order, analytics, dan kondisi integrasi.

Marketplace seperti PasarKita membutuhkan algoritma karena data yang ditampilkan tidak cukup hanya disimpan dan diambil apa adanya. Produk perlu direkomendasikan, diurutkan, dicari, dan dihitung biayanya. Recommender system membantu pengguna memilih item ketika pilihan terlalu banyak \parencite{Resnick1997}. Dari sisi organisasi, penggunaan teknologi informasi juga relevan untuk membantu proses adopsi digital pada level usaha atau organisasi \parencite{Oliveira2011}. Karena itu, PasarKita cocok dijadikan objek analisis penerapan algoritma dasar.

\subsection{Fokus Analisis}

Laporan ini hanya membahas fitur yang sudah memiliki penerapan algoritma jelas di dokumentasi dan kode aplikasi. Fokusnya adalah sebagai berikut:

\begin{itemize}[leftmargin=*]
    \item Algoritma Greedy pada rekomendasi produk, sorting produk, dan fee marketplace.
    \item String Matching/KMP pada pencarian produk buyer dan seller.
    \item String Matching multi-field pada pencarian pesanan.
    \item Potongan kode, pseudocode, dan contoh output dari penerapan algoritma.
\end{itemize}

Laporan ini tidak mengklaim bahwa PasarKita sudah memakai machine learning recommendation. Referensi Epsilon Greedy digunakan sebagai catatan pengembangan, bukan sebagai implementasi saat ini \parencite{Umami2021}.

\subsection{Fitur yang Dianalisis}

Fitur yang dianalisis dipilih karena langsung berhubungan dengan kebutuhan marketplace.

\begin{table}[htbp]
    \caption{Fitur PasarKita yang Dianalisis}
    \label{tab:fitur-dianalisis}
    \small
    \begin{tabularx}{\textwidth}{@{} >{\raggedright\arraybackslash}X >{\raggedright\arraybackslash}X >{\raggedright\arraybackslash}X @{}}
    \toprule
    \textbf{Fitur} & \textbf{Masalah yang Diselesaikan} & \textbf{Algoritma} \\
    \midrule
    Rekomendasi produk & Pembeli butuh produk yang relevan dari banyak pilihan & Greedy scoring \\
    Sorting produk & Produk perlu ditampilkan berdasarkan performa penjualan atau rating & Greedy ranking \\
    Fee marketplace & Total checkout perlu dihitung cepat dan konsisten & Greedy satu langkah \\
    Pencarian produk & Buyer dan seller perlu menemukan produk berdasarkan keyword & KMP/String Matching \\
    Pencarian pesanan & Seller perlu mencari order berdasarkan beberapa field & String Matching multi-field \\
    \bottomrule
    \end{tabularx}
\end{table}

\subsection{Sumber Data Analisis}

Analisis dilakukan dari tiga sumber utama. Pertama, dokumentasi algoritma pada folder \path{docs/algorithms}. Kedua, source code aplikasi yang berkaitan dengan rekomendasi, fee, pencarian produk, dan pencarian pesanan. File utama yang ditinjau adalah modul rekomendasi frontend, utilitas fee, utilitas KMP, service produk, dan service order. Ketiga, referensi yang sudah dikumpulkan dalam file BibTeX \path{src/My Collection.bib}.


\clearpage

% BAB 2 - GAMBARAN FITUR DAN REFERENSI
\section*{BAB 2 \\ GAMBARAN FITUR DAN REFERENSI}
\phantomsection
\addcontentsline{toc}{section}{BAB 2 GAMBARAN FITUR DAN REFERENSI}
\refstepcounter{section}

\subsection{Analisis Jurnal sebagai Dasar Pemilihan Algoritma}

Referensi pada bagian ini tidak hanya dipakai sebagai kutipan teori, tetapi sebagai alasan mengapa algoritma tertentu masuk akal untuk PasarKita. Cara membacanya sederhana: jurnal memberi konteks masalah, lalu konteks itu diterjemahkan menjadi keputusan fitur di aplikasi.

\begingroup
\scriptsize
\renewcommand{\arraystretch}{1.15}
\begin{longtable}{@{} >{\raggedright\arraybackslash}p{0.22\textwidth} >{\raggedright\arraybackslash}p{0.35\textwidth} >{\raggedright\arraybackslash}p{0.35\textwidth} @{}}
    \caption{Analisis Jurnal dan Hubungannya dengan PasarKita}
    \label{tab:analisis-jurnal}\\
    \toprule
    \textbf{Referensi} & \textbf{Isi penting jurnal} & \textbf{Hubungan dengan Aplikasi} \\
    \midrule
    \endfirsthead
    \caption[]{Analisis Jurnal dan Hubungannya dengan PasarKita (lanjutan)}\\
    \toprule
    \textbf{Referensi} & \textbf{Isi penting jurnal} & \textbf{Hubungan dengan Aplikasi} \\
    \midrule
    \endhead
    \bottomrule
    \endfoot
    \textcite{Oliveira2011} & Membahas model adopsi teknologi pada level organisasi, termasuk kerangka DOI dan TOE. Intinya, teknologi dipakai ketika bisa membantu organisasi bekerja lebih efisien dan menyesuaikan diri dengan lingkungan bisnis. & PasarKita dibuat sebagai kanal digital untuk UMKM. Karena itu, fitur marketplace tidak cukup hanya menampilkan data; aplikasi perlu membantu proses jual-beli berjalan lebih cepat, mudah dicari, dan mudah dipantau. \\
    \addlinespace
    \textcite{Resnick1997} & Menjelaskan bahwa recommender system membantu user memilih item ketika pilihan terlalu banyak. Masalah utamanya bukan kekurangan data, tetapi terlalu banyak alternatif yang harus dipilih user. & Produk UMKM dalam marketplace bisa bertambah banyak. PasarKita memakai rekomendasi berbasis skor agar buyer tidak harus menelusuri semua produk satu per satu. \\
    \addlinespace
    \textcite{Tanujaya2026} & Membandingkan Greedy dan Dynamic Programming pada optimasi diskon keranjang e-commerce. Greedy ditunjukkan sebagai pendekatan yang praktis dan cepat, walaupun tidak selalu menjamin hasil optimal global seperti pendekatan yang lebih berat. & Fitur rekomendasi, sorting, dan fee PasarKita membutuhkan respons cepat. Karena targetnya ranking dan perhitungan sederhana, Greedy lebih cocok daripada algoritma optimasi yang terlalu kompleks. \\
    \addlinespace
    \textcite{Marbun2024} & Menerapkan Knuth-Morris-Pratt pada e-katalog perpustakaan. Fokusnya adalah mencocokkan kata kunci dengan data katalog agar pencarian item lebih efisien. & PasarKita memiliki kasus yang mirip dengan e-katalog: user mengetik keyword, lalu sistem mencari produk yang namanya mengandung keyword tersebut. Karena itu, KMP sesuai untuk pencarian produk. \\
    \addlinespace
    \textcite{Al-howaide} & Membahas algoritma string matching dan perbandingan cara pencocokan string. Poin pentingnya adalah setiap algoritma pencarian memiliki cara mengurangi perbandingan karakter yang tidak perlu. & Pencarian produk dan pesanan tidak boleh terasa lambat ketika data bertambah. String Matching dipakai agar keyword bisa dicocokkan ke nama produk, ID transaksi, atau nomor tracking secara terarah. \\
    \addlinespace
    \textcite{Crochemore2013} & Menjelaskan pattern matching sebagai masalah dasar dalam pemrosesan string: mencari kemunculan pattern di dalam teks. & Fitur pencarian PasarKita pada dasarnya adalah pattern matching. Keyword adalah pattern, sedangkan nama produk/order field adalah teks yang dicari. \\
\end{longtable}
\endgroup

\subsection{Alasan Khusus Aplikasi Menggunakan Algoritma Ini}

Tabel berikut bisa dipakai sebagai ringkasan utama saat presentasi. Isinya langsung menghubungkan masalah aplikasi, alasan pemilihan algoritma, dan output yang terlihat.

\begingroup
\scriptsize
\renewcommand{\arraystretch}{1.2}
\setlength{\tabcolsep}{3pt}
\begin{longtable}{@{} >{\raggedright\arraybackslash}p{0.18\textwidth} >{\raggedright\arraybackslash}p{0.25\textwidth} >{\raggedright\arraybackslash}p{0.30\textwidth} >{\raggedright\arraybackslash}p{0.20\textwidth} @{}}
    \caption{Keputusan Algoritma pada Fitur PasarKita}
    \label{tab:keputusan-algoritma}\\
    \toprule
    \textbf{Algoritma} & \textbf{Masalah di aplikasi} & \textbf{Mengapa cocok} & \textbf{Penerapan pada Aplikasi} \\
    \midrule
    \endfirsthead
    \caption[]{Keputusan Algoritma pada Fitur PasarKita (lanjutan)}\\
    \toprule
    \textbf{Algoritma} & \textbf{Masalah di aplikasi} & \textbf{Mengapa cocok} & \textbf{Penerapan pada Aplikasi} \\
    \midrule
    \endhead
    \bottomrule
    \endfoot
    Greedy scoring & Buyer butuh rekomendasi produk yang relevan tanpa proses berat. & Mengambil produk dengan skor terbaik berdasarkan kategori yang sering dilihat, jumlah terjual, dan rating. Cocok untuk rekomendasi awal karena cepat dan mudah dijelaskan. & Daftar produk rekomendasi: stok tersedia, belum dilihat, dan kategori/skornya relevan. \\
    \addlinespace
    Greedy ranking & Produk perlu diurutkan berdasarkan terlaris atau rating tertinggi. & Setiap produk punya skor lokal. Sistem cukup mengurutkan skor tersebut tanpa mencari kombinasi global. & Tabel produk terurut dari sold units atau rating terbesar. \\
    \addlinespace
    Greedy satu langkah & Checkout perlu menghitung fee marketplace dengan konsisten. & Fee selalu 2\% dari subtotal, sehingga cukup dihitung langsung dalam satu langkah. & Subtotal, fee marketplace, dan total pembayaran. \\
    \addlinespace
    KMP & Buyer dan seller perlu mencari produk berdasarkan keyword. & KMP memakai prefix table sehingga ketika terjadi mismatch, pencarian tidak selalu diulang dari awal. Ini cocok untuk exact substring search pada nama produk. & Keyword \texttt{sep} menghasilkan produk seperti \texttt{Sepatu Running Lokal}. \\
    \addlinespace
    String Matching multi-field & Seller perlu mencari order melalui order ID, nama produk, transaction ID, atau tracking ID. & Masalahnya bukan hanya mencari satu nama produk, tetapi mencocokkan keyword ke beberapa field. Pendekatan string matching membuat pencarian lebih fleksibel. & Keyword transaksi atau tracking menampilkan order yang cocok. \\
\end{longtable}
\endgroup

Dari tabel tersebut, Greedy dipilih ketika PasarKita perlu mengambil keputusan cepat berdasarkan skor lokal. KMP dan String Matching dipilih ketika fitur membutuhkan pencocokan keyword pada teks. Referensi \textcite{Umami2021} tidak dipakai untuk mengklaim implementasi saat ini, tetapi menjadi catatan pengembangan jika rekomendasi PasarKita ingin dibuat lebih adaptif dengan pola eksplorasi dan eksploitasi.

\subsection{Alur Fitur yang Dianalisis}

Secara umum, fitur yang dianalisis mengikuti alur input, proses algoritma, dan output. Rekomendasi produk menerima input berupa daftar produk dan histori produk yang pernah dilihat. Sistem menghitung skor setiap produk, mengurutkan produk berdasarkan skor, lalu mengambil produk teratas. Sorting produk bekerja dengan pola serupa, tetapi skor utamanya berasal dari unit terjual atau rating. Fee marketplace menerima subtotal dan langsung menghasilkan fee serta total pembayaran.

Pada pencarian, input utama adalah keyword. Keyword tersebut dicocokkan dengan nama produk atau field order. Untuk pencarian produk, PasarKita menggunakan KMP yang bersifat case-insensitive. Untuk pencarian pesanan, keyword dapat dicari pada order ID, nama produk, transaction ID, dan tracking ID.

\subsection{Ringkasan Input, Proses, dan Output}

\begin{table}[htbp]
    \caption{Ringkasan Input, Proses, dan Output Algoritma}
    \label{tab:ipo-algoritma}
    \small
    \begin{tabularx}{\textwidth}{@{} >{\raggedright\arraybackslash}X >{\raggedright\arraybackslash}X >{\raggedright\arraybackslash}X >{\raggedright\arraybackslash}X @{}}
    \toprule
    \textbf{Fitur} & \textbf{Input} & \textbf{Proses} & \textbf{Output} \\
    \midrule
    Rekomendasi produk & Produk dan histori lihat & Hitung skor lokal, sort, ambil top-N & Produk rekomendasi \\
    Sorting produk & Data produk, sold units, rating & Ranking berdasarkan skor utama & Produk terurut \\
    Fee marketplace & Subtotal checkout & Hitung 2\% dari subtotal & Fee dan total \\
    Pencarian produk & Keyword dan nama produk & KMP case-insensitive & Produk yang cocok \\
    Pencarian pesanan & Keyword dan field order & Matching beberapa field & Order yang cocok \\
    \bottomrule
    \end{tabularx}
\end{table}


\clearpage

% BAB 3 - ANALISIS ALGORITMA DAN OUTPUT
\section*{BAB 3 \\ ANALISIS ALGORITMA DAN OUTPUT}
\phantomsection
\addcontentsline{toc}{section}{BAB 3 ANALISIS ALGORITMA DAN OUTPUT}
\refstepcounter{section}

\subsection{Deskripsi Solusi}

CleanConnect dirancang sebagai sebuah platform digital berbasis aplikasi seluler (mobile application) yang berfungsi sebagai penghubung antara pengguna jasa (konsumen) dengan penyedia jasa kebersihan rumah tangga (mitra petugas kebersihan) melalui mekanisme pemesanan berbasis lokasi. Platform ini dikembangkan dengan pendekatan model bisnis on-demand service, yang mengadaptasi prinsip-prinsip dari platform layanan berbasis lokasi yang telah berhasil diterapkan pada sektor lain, namun disesuaikan dengan karakteristik dan kebutuhan spesifik layanan kebersihan rumah tangga \cite{Muangmee2021,An2023}.

Secara umum, CleanConnect terdiri atas dua antarmuka utama, yaitu antarmuka pengguna (customer application) dan antarmuka mitra petugas (partner application), yang keduanya terhubung melalui sistem backend terpusat. Sistem ini mengelola data pemesanan, lokasi, estimasi harga, riwayat transaksi, serta rekam jejak layanan. Berdasarkan kajian masalah pada Bab II, solusi CleanConnect diarahkan untuk menjawab tiga isu utama, yaitu kesulitan menemukan penyedia jasa terpercaya, kurangnya transparansi harga, dan ketidakpastian waktu kedatangan petugas \cite{Hawlitschek2016,Li2020}.

Fitur utama yang diusulkan meliputi booking cleaning service, tracking petugas berbasis lokasi/GPS, rating dan review, estimasi biaya otomatis, serta jadwal rutin mingguan atau bulanan. Fitur booking cleaning service memungkinkan pengguna melakukan pemesanan layanan secara langsung melalui aplikasi tanpa harus mencari penyedia jasa secara manual. Fitur tracking petugas membantu pengguna memantau posisi petugas secara real-time, sedangkan fitur rating dan review menjadi mekanisme reputasi untuk menilai kualitas layanan. Fitur estimasi biaya otomatis memberikan gambaran biaya sebelum pemesanan dikonfirmasi, dan fitur jadwal rutin membantu pengguna mengatur layanan berkala tanpa melakukan pemesanan ulang setiap kali membutuhkan layanan \cite{Pena-garcia2024,Id2018}.

\subsection{Teknologi yang Digunakan}

Teknologi yang digunakan dalam rancangan CleanConnect dapat dikelompokkan ke dalam beberapa komponen utama. Komponen pertama adalah aplikasi seluler yang menjadi media interaksi antara pengguna dan sistem. Aplikasi ini digunakan untuk melakukan registrasi, memilih layanan, menentukan lokasi, mengatur jadwal, melihat estimasi biaya, memantau status petugas, dan memberikan rating setelah layanan selesai \cite{An2023,Nalendra2024}.

Komponen kedua adalah layanan lokasi atau GPS yang digunakan untuk mendukung fitur tracking petugas. Melalui layanan ini, posisi petugas dapat dikirimkan secara periodik ke sistem backend dan ditampilkan pada peta digital di aplikasi pengguna. Informasi lokasi tersebut juga dapat digunakan untuk menampilkan estimasi waktu kedatangan petugas.

Komponen ketiga adalah sistem backend dan basis data. Backend berfungsi mengelola proses pemesanan, pencocokan petugas, penghitungan estimasi biaya, pengelolaan profil pengguna dan mitra petugas, penyimpanan riwayat transaksi, serta pengelolaan rating dan review. Basis data digunakan untuk menyimpan informasi pengguna, mitra petugas, lokasi layanan, jadwal, status pemesanan, dan ulasan layanan \cite{Pena-garcia2024,Id2018}.

Komponen keempat adalah sistem notifikasi yang digunakan untuk memberi informasi kepada pengguna dan petugas mengenai perubahan status pesanan, konfirmasi layanan, estimasi kedatangan, serta pengingat jadwal rutin. Dengan kombinasi komponen tersebut, CleanConnect diharapkan dapat menghadirkan proses pemesanan layanan kebersihan yang lebih praktis, transparan, dan mudah dipantau.


% DAFTAR PUSTAKA
\clearpage
\section*{DAFTAR PUSTAKA}
\phantomsection
\addcontentsline{toc}{section}{DAFTAR PUSTAKA}

\printbibliography[heading=none]

\end{document}
